\section{Background and Motivation}
\label{sec:intro:background}

\subsection{The Grasping Paradigm and Its Limits}
\label{sec:intro:grasping}

Humans engage naturally in multiple forms of dexterous manipulation that involve grasping, pushing, poking, rolling, and tossing objects \cite{bullock2011classifying}.
This highlights a significant gap between human and robotic manipulation capabilities, as evidenced by the disproportionate research attention devoted to prehensile manipulation, or \textit{grasping} \cite{billard2019trends}.
Grasping is computationally attractive primarily because, once an object is under form or force closure, it generally does not need to be tracked over time and uncertainty on its state is reduced.

Despite these advantages, grasping is limited in four ways.
First, it is constrained by the \textit{reachability} of the robot arm---objects outside the manipulator's workspace cannot be grasped regardless of the planner's sophistication.
Second, it depends on \textit{end-effector design}---gripper geometry determines which object shapes, sizes, and orientations are graspable.
Third, it depends on the \textit{physical properties} of the object---soft, fragile, or deformable objects may not survive grasp forces.
Fourth, it is sensitive to the \textit{accuracy of the perception system}---small pose estimation errors propagate directly into grasp failure.

Beyond these constraints, extending prehensile capability to objects not rigidly attached to the end-effector introduces the challenge of coupled dynamics between robot and object motion \cite{yin2021modeling,ichnowski2022gomp}.
Tasks like liquid transport, high-payload handling, and manipulation after joint failures require reasoning about torques, inertial forces, and contact mechanics---phenomena that geometric planners cannot represent.

\subsection{Two Paradigms, Two Limitations}
\label{sec:intro:paradigms}

Recent advances in robotic manipulation have proceeded along two largely separate paths with complementary strengths and weaknesses.

\header{Traditional control-based methods}
Classical control and optimization approaches successfully incorporate robot dynamics through analytical models and constrained optimization, but struggle to scale beyond carefully engineered environments due to modeling complexity \cite{feng2014optimization,gamba2023springy}.
These methods provide strong theoretical guarantees---stability certificates, constraint satisfaction bounds, optimality conditions---and have been deployed successfully in industrial settings.
Their limitation is scalability: accurate analytical models are hard to construct for complex robot-object interactions, and the engineering effort to extend these systems to new tasks is substantial.

\header{Learning-based approaches}
Data-driven methods, including diffusion policies and foundation models for manipulation, have demonstrated impressive real-world generalization through imitation learning \cite{chi2023diffusion,team2024octo}.
These approaches reduce task-specific engineering burden and improve naturally over deployment lifespans.
Their limitation is depth: current methods typically operate through position commands, focus on prehensile manipulation, and treat dynamics as an implicit property of demonstrated trajectories rather than an explicit constraint to satisfy.
A policy that has never seen a joint failure, a heavy payload, or a liquid-filled container will not generalize to them.

The fundamental disconnect---between dynamics-aware but brittle traditional approaches and generalizable but dynamics-limited learning methods---is a critical gap in advancing manipulation capabilities \cite{ruggiero2018nonprehensile}.

\subsection{Four Recurring Challenges}
\label{sec:intro:challenges}

Across both paradigms, four challenges recur whenever dynamics enter the manipulation problem.
These challenges are not incidental; they reflect deep structural tensions in the problem.

\textit{Modeling burden}: constructing or learning an accurate dynamics model for complex interactions is expensive.
Simulation-in-the-loop approaches provide accuracy but require repeated forward propagation at planning time; learned models offer speed but are limited to the dynamics they have seen.

\textit{Dimensionality}: incorporating higher-order state---velocity, acceleration, torque---alongside geometric constraints exponentially expands the search and optimization space.
Kinodynamic planning in joint-velocity-acceleration space is fundamentally harder than geometric path planning, and the difficulty grows with system complexity.

\textit{Integration and pipeline brittleness}: combining a dynamics model, a planner, a trajectory optimizer, and safety checks into a working system introduces brittleness at every interface.
Each component requires careful tuning, and errors propagate across components in ways that are hard to debug or predict.

\textit{Usability}: traditional dynamics-constrained methods require task-specific engineering---developing bespoke sampling strategies, designing cost functions, determining which phenomena to model---that limits deployment to controlled environments.

\cref{ch:framework} provides a detailed treatment of these challenges and organizes the landscape of dynamics-constrained motion generation around them.
They recur throughout this thesis as a lens for evaluating each contribution.
